\documentclass[10pt,a4paper]{amsart}
\usepackage{amsmath,amssymb,amsthm,graphicx}
\usepackage{verbatim}
\newenvironment{reference}{\comment}{\endcomment}
%\newenvironment{reference}{}{}
\newenvironment{slogan}{\comment}{\endcomment}
\newenvironment{history}{\comment}{\endcomment}

% For commutative diagrams we use Xy-pic
\usepackage[all]{xy}

% We use 2cell for 2-commutative diagrams.
\xyoption{2cell}
% Theorem environments.
%
\theoremstyle{plain}
\newtheorem{theorem}[subsection]{Theorem}
\newtheorem{proposition}[subsection]{Proposition}
\newtheorem{lemma}[subsection]{Lemma}

\theoremstyle{definition}
\newtheorem{definition}[subsection]{Definition}
\newtheorem{example}[subsection]{Example}
\newtheorem{exercise}[subsection]{Exercise}
\newtheorem{situation}[subsection]{Situation}

\theoremstyle{remark}
\newtheorem{remark}[subsection]{Remark}
\newtheorem{remarks}[subsection]{Remarks}

\numberwithin{equation}{subsection}

% Macros
%
\def\lim{\mathop{\mathrm{lim}}\nolimits}
\def\colim{\mathop{\mathrm{colim}}\nolimits}
\def\Spec{\mathop{\mathrm{Spec}}}
\def\Hom{\mathop{\mathrm{Hom}}\nolimits}
\def\Ext{\mathop{\mathrm{Ext}}\nolimits}
\def\SheafHom{\mathop{\mathcal{H}\!\mathit{om}}\nolimits}
\def\SheafExt{\mathop{\mathcal{E}\!\mathit{xt}}\nolimits}
\def\Sch{\mathit{Sch}}
\def\Mor{\mathop{\mathrm{Mor}}\nolimits}
\def\Ob{\mathop{\mathrm{Ob}}\nolimits}
\def\Sh{\mathop{\mathit{Sh}}\nolimits}
\def\NL{\mathop{N\!L}\nolimits}
\def\CH{\mathop{\mathrm{CH}}\nolimits}
\def\proetale{{pro\text{-}\acute{e}tale}}
\def\etale{{\acute{e}tale}}
\def\QCoh{\mathit{QCoh}}
\def\Ker{\mathop{\mathrm{Ker}}}
\def\Im{\mathop{\mathrm{Im}}}
\def\Coker{\mathop{\mathrm{Coker}}}
\def\Coim{\mathop{\mathrm{Coim}}}

% Boxtimes
%
\DeclareMathSymbol{\boxtimes}{\mathbin}{AMSa}{"02}

%
% Macros for moduli stacks/spaces
%
\def\QCohstack{\mathcal{QC}\!\mathit{oh}}
\def\Cohstack{\mathcal{C}\!\mathit{oh}}
\def\Spacesstack{\mathcal{S}\!\mathit{paces}}
\def\Quotfunctor{\mathrm{Quot}}
\def\Hilbfunctor{\mathrm{Hilb}}
\def\Curvesstack{\mathcal{C}\!\mathit{urves}}
\def\Polarizedstack{\mathcal{P}\!\mathit{olarized}}
\def\Complexesstack{\mathcal{C}\!\mathit{omplexes}}
% \Pic is the operator that assigns to X its picard group, usage \Pic(X)
% \Picardstack_{X/B} denotes the Picard stack of X over B
% \Picardfunctor_{X/B} denotes the Picard functor of X over B
\def\Pic{\mathop{\mathrm{Pic}}\nolimits}
\def\Picardstack{\mathcal{P}\!\mathit{ic}}
\def\Picardfunctor{\mathrm{Pic}}
\def\Deformationcategory{\mathcal{D}\!\mathit{ef}}
\usepackage{cleveref}
\crefname{originallemma}{original lemma}{original lemmas}
\Crefname{originallemma}{Original lemma}{Original lemmas}
\makeatletter
\newlabel{lemma-criterion-baer}{{05NU}{source}{}{sourcepointer-1}{}}
\newlabel{lemma-criterion-baer@cref}{{[originallemma][0][]05NU}{[1][0][]source}{}{}{}}
\makeatother
\crefname{originalproposition}{original proposition}{original propositions}
\Crefname{originalproposition}{Original proposition}{Original propositions}
\makeatletter
\newlabel{proposition-modules-are-small}{{05NT}{source}{}{sourcepointer-2}{}}
\newlabel{proposition-modules-are-small@cref}{{[originalproposition][0][]05NT}{[1][0][]source}{}{}{}}
\makeatother
\crefname{originallemma}{original lemma}{original lemmas}
\Crefname{originallemma}{Original lemma}{Original lemmas}
\makeatletter
\newlabel{lemma-construction}{{05NW}{source}{}{sourcepointer-3}{}}
\newlabel{lemma-construction@cref}{{[originallemma][0][]05NW}{[1][0][]source}{}{}{}}
\makeatother
\title{Baer's criterion produces functorial injective embeddings for modules}
\author{The Stacks Project Authors}
\date{stacks-project:injectives.tex:local-snapshot-20261012; GNU FDL 1.2 or later}
\begin{document}
\maketitle
\noindent\textit{Editorial note.} All delivered mathematical text is the author's own, copied from the pinned Stacks Project chapter file sites-cohomology.tex; no statement or proof step was rewritten, reordered or removed.\par\medskip
\noindent\textit{Editorial note.} The proof is the author's complete argument; the extension step is displayed with the relevant inclusions and the cofinality hypothesis is used explicitly, so no step is left to the reader.\par\medskip
\noindent\textit{Editorial note.} Every labelled reference occurring inside the counted spans is carried in assembly.external\_references as a non-load-bearing navigation pointer (same-chapter ones at their definition site, cross-chapter ones at the invocation inside the counted span), so no load-bearing mathematics is replaced by a stub. No bibliography citation occurs inside the counted spans.\par\medskip
\noindent\textit{Editorial note.} The same-name check for this card scanned every staged Stacks chapter of this repository -- all sources/stacks-*/pinned-original-source.tex files, the card's own chapter excluded -- and the labels found are recorded in dedup.same\_name\_cross\_chapter on the card.\par\medskip
\noindent\textit{Editorial note.} The preamble is the chapter's own retained minimal preamble (sources/stacks-sites-cohomology/retained-preamble.tex), a verbatim excerpt of the native preamble: the Xy-pic setup and the native theorem declarations with the mathematical macros.\par\medskip
\noindent\textit{Allowed standard prerequisites (editorial).}\begin{itemize}
\item Baer's criterion. Same chapter, Lemma lemma-criterion-baer. Supplies that it suffices to extend maps from ideals of the ring, which is how the proof is organized.
\item The transfinite construction of the embedding. Same chapter, the construction of the modules attached to an ordinal and the functor M used in it. Supplies the objects whose injectivity is being proved, together with the functoriality of the construction and the inclusions between successive stages.
\end{itemize}\medskip
\ifdefined\hypertarget\hypertarget{sourcepointer-1}{}\fi
\noindent\textit{Source reference.} Original lemma 05NU: sources/stacks-injectives/pinned-original-source.tex, source lines 437--437. Navigation pointer invoked at proof line 576; the counted span contains the pointer only and the cited result is not reproduced inside it, so the pointer is registered as a non-load-bearing navigation stub in the established form for Stacks material. No bibliography citation is involved.\par\medskip
\ifdefined\hypertarget\hypertarget{sourcepointer-2}{}\fi
\noindent\textit{Source reference.} Original proposition 05NT: sources/stacks-injectives/pinned-original-source.tex, source lines 387--387. Navigation pointer invoked at proof line 586; the counted span contains the pointer only and the cited result is not reproduced inside it, so the pointer is registered as a non-load-bearing navigation stub in the established form for Stacks material. No bibliography citation is involved.\par\medskip
\ifdefined\hypertarget\hypertarget{sourcepointer-3}{}\fi
\noindent\textit{Source reference.} Original lemma 05NW: sources/stacks-injectives/pinned-original-source.tex, source lines 500--500. Navigation pointer invoked at proof line 599; the counted span contains the pointer only and the cited result is not reproduced inside it, so the pointer is registered as a non-load-bearing navigation stub in the established form for Stacks material. No bibliography citation is involved.\par\medskip
\begin{theorem}
\label{theorem-baer-grothendieck}
Let $\kappa$ be the cardinality of the set of ideals in $R$, and
let $\alpha$ be an ordinal whose cofinality is greater than
$\kappa$. Then $\mathbf{M}_\alpha(N)$ is an injective $R$-module,
and $N \to \mathbf{M}_\alpha(N)$ is a functorial injective embedding.
\end{theorem}

\begin{proof}
By Baer's criterion
Lemma \ref{lemma-criterion-baer},
it suffices to show that if $\mathfrak{a} \subset R$ is an ideal, then
any map $f : \mathfrak{a} \to \mathbf{M}_\alpha(N)$ extends to
$R \to \mathbf{M}_\alpha(N)$. However, we know since $\alpha$ is a limit
ordinal that
$$
\mathbf{M}_{\alpha}(N) =
\colim_{\beta < \alpha} \mathbf{M}_{\beta}(N),
$$
so by
Proposition \ref{proposition-modules-are-small},
we find that
$$
\Hom_R(\mathfrak{a}, \mathbf{M}_{\alpha}(N)) =
\colim_{\beta < \alpha} \Hom_R(\mathfrak a, \mathbf{M}_{\beta}(N)).
$$
This means in particular that there is some $\beta' < \alpha$
such that $f$ factors through the submodule $\mathbf{M}_{\beta'}(N)$, as
$$
f : \mathfrak{a} \to \mathbf{M}_{\beta'}(N) \to
\mathbf{M}_{\alpha}(N).
$$
However, by the fundamental property of the functor $\mathbf{M}$,
see Lemma \ref{lemma-construction} part (3),
we know that the map $\mathfrak{a} \to \mathbf{M}_{\beta'}(N)$
can be extended to
$$
R \to \mathbf{M}( \mathbf{M}_{\beta'}(N)) =
\mathbf{M}_{\beta' + 1}(N),
$$
and the last object imbeds in $\mathbf{M}_{\alpha}(N)$ (as
$\beta' + 1 < \alpha$ since $\alpha$ is a limit ordinal).
In particular, $f$ can be extended to $\mathbf{M}_{\alpha}(N)$.
\end{proof}

\end{document}
